\textbf{Phần 1. Trắc nghiệm lựa chọn}
\setcounterpageref{ex}{0}

\begin{ex} %[5D1V7]
Tính giá trị của biểu thức $P = 3\dfrac{2}{5} + 4\dfrac{3}{7} + 2\dfrac{3}{5} + 1\dfrac{4}{7}$ bằng cách thuận tiện nhất:
\choice
{$11$}
{\True $12$}
{$10$}
{$13$}
\loigiai{
Nhóm các hỗn số có cùng mẫu số ở phần phân số:
\begin{align*}
P &= \left(3\dfrac{2}{5} + 2\dfrac{3}{5}\right) + \left(4\dfrac{3}{7} + 1\dfrac{4}{7}\right) \\
&= \left(3 + 2 + \dfrac{2}{5} + \dfrac{3}{5}\right) + \left(4 + 1 + \dfrac{3}{7} + \dfrac{4}{7}\right) \\
&= (5 + 1) + (5 + 1) = 6 + 6 = 12.
\end{align*}
\textbf{Đáp án của câu hỏi là B}
}
\end{ex}

\begin{ex} %[5D1V7]
Tìm $x$ biết: $x - 2\dfrac{1}{4} = 3\dfrac{1}{2}$.
\choice
{$x = 5\dfrac{1}{4}$}
{\True $x = 5\dfrac{3}{4}$}
{$x = 6\dfrac{1}{4}$}
{$x = 5\dfrac{1}{2}$}
\loigiai{
Ta có:
\[
x = 3\dfrac{1}{2} + 2\dfrac{1}{4} = \left(3 + \dfrac{2}{4}\right) + \left(2 + \dfrac{1}{4}\right) = 5 + \dfrac{3}{4} = 5\dfrac{3}{4}.
\]
\textbf{Đáp án của câu hỏi là B}
}
\end{ex}

\begin{ex} %[5D1V7]
Chuyển hỗn số $4\dfrac{3}{8}$ thành phân số rồi nhân với $\dfrac{4}{7}$, kết quả rút gọn là
\choice
{\True $\dfrac{5}{2}$}
{$2$}
{$\dfrac{7}{2}$}
{$3$}
\loigiai{
Chuyển hỗn số thành phân số:
\[
4\dfrac{3}{8} = \dfrac{4 \times 8 + 3}{8} = \dfrac{35}{8}.
\]
Thực hiện phép nhân:
\[
\dfrac{35}{8} \times \dfrac{4}{7} = \dfrac{35 \times 4}{8 \times 7} = \dfrac{5 \times 1}{2 \times 1} = \dfrac{5}{2}.
\]
\textbf{Đáp án của câu hỏi là A}
}
\end{ex}

\begin{ex} %[5D1V7]
Một tấm vải dài $10\text{ m}$. Người thợ may cắt ra $3\dfrac{1}{2}\text{ m}$ để may áo và $2\dfrac{1}{4}\text{ m}$ để may quần. Số mét vải còn lại là
\choice
{$4\dfrac{1}{4}\text{ m}$}
{\True $4\dfrac{1}{4}\text{ m}$}
{$5\dfrac{3}{4}\text{ m}$}
{$4\dfrac{3}{4}\text{ m}$}
\loigiai{
Tổng số mét vải đã dùng là:
\[
3\dfrac{1}{2} + 2\dfrac{1}{4} = 3\dfrac{2}{4} + 2\dfrac{1}{4} = 5\dfrac{3}{4}\text{ (m)}.
\]
Số mét vải còn lại là:
\[
10 - 5\dfrac{3}{4} = 9\dfrac{4}{4} - 5\dfrac{3}{4} = 4\dfrac{1}{4}\text{ (m)}.
\]
\textbf{Đáp án của câu hỏi là A}
}
\end{ex}

\begin{ex} %[5D1V7]
Trong các hỗn số sau, hỗn số nào có giá trị lớn nhất?
\choice
{$3\dfrac{4}{5}$}
{$3\dfrac{7}{10}$}
{\True $4\dfrac{1}{8}$}
{$3\dfrac{9}{10}$}
\loigiai{
So sánh phần nguyên của các hỗn số:
Các hỗn số $3\dfrac{4}{5}$, $3\dfrac{7}{10}$, $3\dfrac{9}{10}$ đều có phần nguyên là $3$.\\
Hỗn số $4\dfrac{1}{8}$ có phần nguyên là $4$.\\
Vì $4 > 3$ nên hỗn số lớn nhất là $4\dfrac{1}{8}$.\\
\textbf{Đáp án của câu hỏi là C}
}
\end{ex}

\begin{ex} %[5D1V7]
Một hình chữ nhật có chiều dài là $4\dfrac{1}{2}\text{ m}$, chiều rộng kém chiều dài $1\dfrac{1}{4}\text{ m}$. Chu vi của hình chữ nhật đó là
\choice
{$14\text{ m}$}
{\True $15\dfrac{1}{2}\text{ m}$}
{$7\dfrac{3}{4}\text{ m}$}
{$16\text{ m}$}
\loigiai{
Chiều rộng hình chữ nhật là:
\[
4\dfrac{1}{2} - 1\dfrac{1}{4} = 4\dfrac{2}{4} - 1\dfrac{1}{4} = 3\dfrac{1}{4}\text{ (m)}.
\]
Nửa chu vi hình chữ nhật là:
\[
4\dfrac{1}{2} + 3\dfrac{1}{4} = 4\dfrac{2}{4} + 3\dfrac{1}{4} = 7\dfrac{3}{4}\text{ (m)}.
\]
Chu vi hình chữ nhật là:
\[
7\dfrac{3}{4} \times 2 = \dfrac{31}{4} \times 2 = \dfrac{31}{2} = 15\dfrac{1}{2}\text{ (m)}.
\]
\textbf{Đáp án của câu hỏi là B}
}
\end{ex}

\begin{ex} %[5D1V7]
Tìm số tự nhiên $n$ biết: $3\dfrac{2}{7} < n < 5\dfrac{1}{3}$.
\choice
{$n = 4$}
{$n = 5$}
{\True $n \in \{4, 5\}$}
{$n \in \{3, 4, 5\}$}
\loigiai{
Ta thấy: $3 < 3\dfrac{2}{7} < 4 < 5 < 5\dfrac{1}{3} < 6$.\\
Các số tự nhiên nằm giữa $3\dfrac{2}{7}$ và $5\dfrac{1}{3}$ là $4$ và $5$.\\
\textbf{Đáp án của câu hỏi là C}
}
\end{ex}

\begin{ex} %[5D1V7]
Một xe ô tô giờ thứ nhất chạy được $45\dfrac{1}{2}\text{ km}$, giờ thứ hai chạy được $48\dfrac{3}{4}\text{ km}$, giờ thứ ba chạy được $46\dfrac{1}{4}\text{ km}$. Hỏi trung bình mỗi giờ xe ô tô chạy được bao nhiêu ki-lô-mét?
\choice
{$46\text{ km}$}
{\True $46\dfrac{5}{6}\text{ km}$}
{$47\text{ km}$}
{$46\dfrac{1}{2}\text{ km}$}
\loigiai{
Tổng quãng đường xe chạy trong $3$ giờ là:
\[
45\dfrac{1}{2} + 48\dfrac{3}{4} + 46\dfrac{1}{4} = 45\dfrac{1}{2} + \left(48\dfrac{3}{4} + 46\dfrac{1}{4}\right) = 45\dfrac{1}{2} + 95 = 140\dfrac{1}{2} = \dfrac{281}{2}\text{ (km)}.
\]
Trung bình mỗi giờ xe chạy được:
\[
\dfrac{281}{2} : 3 = \dfrac{281}{6} = 46\dfrac{5}{6}\text{ (km)}.
\]
\textbf{Đáp án của câu hỏi là B}
}
\end{ex}

\begin{ex} %[5D1V7]
Chuyển phân số $\dfrac{107}{12}$ thành hỗn số ta được
\choice
{$8\dfrac{1}{12}$}
{\True $8\dfrac{11}{12}$}
{$9\dfrac{1}{12}$}
{$8\dfrac{5}{12}$}
\loigiai{
Thực hiện phép chia: $107 : 12 = 8$ (dư $11$).\\
Vậy $\dfrac{107}{12} = 8\dfrac{11}{12}$.\\
\textbf{Đáp án của câu hỏi là B}
}
\end{ex}

\begin{ex} %[5D1V7]
Tính giá trị biểu thức $Q = 6\dfrac{3}{8} - 2\dfrac{1}{4} + 1\dfrac{5}{8}$:
\choice
{$5$}
{\True $5\dfrac{3}{4}$}
{$5\dfrac{1}{2}$}
{$6$}
\loigiai{
Nhóm các hỗn số có cùng mẫu số:
\[
Q = \left(6\dfrac{3}{8} + 1\dfrac{5}{8}\right) - 2\dfrac{1}{4} = \left(7 + \dfrac{8}{8}\right) - 2\dfrac{1}{4} = 8 - 2\dfrac{1}{4} = 7\dfrac{4}{4} - 2\dfrac{1}{4} = 5\dfrac{3}{4}.
\]
\textbf{Đáp án của câu hỏi là B}
}
\end{ex}

\begin{ex} %[5D1V7]
Có ba túi đường: túi thứ nhất nặng $1\dfrac{3}{4}\text{ kg}$, túi thứ hai nặng hơn túi thứ nhất $\dfrac{1}{2}\text{ kg}$, túi thứ ba nặng bằng cả hai túi kia cộng lại. Cả ba túi đường nặng tất cả bao nhiêu ki-lô-gam?
\choice
{$4\text{ kg}$}
{\True $8\text{ kg}$}
{$6\text{ kg}$}
{$7\dfrac{1}{2}\text{ kg}$}
\loigiai{
Túi thứ hai nặng:
\[
1\dfrac{3}{4} + \dfrac{1}{2} = 1\dfrac{3}{4} + \dfrac{2}{4} = 1\dfrac{5}{4} = 2\dfrac{1}{4}\text{ (kg)}.
\]
Túi thứ nhất và túi thứ hai nặng:
\[
1\dfrac{3}{4} + 2\dfrac{1}{4} = 4\text{ (kg)}.
\]
Túi thứ ba nặng bằng tổng hai túi đầu nên cũng nặng $4\text{ kg}$.\\
Cả ba túi nặng: $4 + 4 = 8\text{ (kg)}$.\\
\textbf{Đáp án của câu hỏi là B}
}
\end{ex}

\begin{ex} %[5D1V7]
Tìm phân số $\dfrac{x}{y}$ biết rằng khi cộng $\dfrac{x}{y}$ với $2\dfrac{1}{3}$ thì được kết quả là $5$:
\choice
{\True $2\dfrac{2}{3}$}
{$3\dfrac{1}{3}$}
{$2\dfrac{1}{3}$}
{$3\dfrac{2}{3}$}
\loigiai{
Ta có:
\[
\dfrac{x}{y} = 5 - 2\dfrac{1}{3} = 4\dfrac{3}{3} - 2\dfrac{1}{3} = 2\dfrac{2}{3}.
\]
\textbf{Đáp án của câu hỏi là A}
}
\end{ex}

\textbf{Phần 2. Trắc nghiệm đúng sai}
\setcounterpageref{ex}{0}

\begin{ex} %[5D1V7]
Cho các hỗn số: $A = 5\dfrac{2}{7}$, $B = 3\dfrac{5}{7}$ và $C = 2\dfrac{3}{14}$. Xét tính đúng sai của các khẳng định sau:
\choiceTF
{\True Tổng $A + B$ là một số tự nhiên bằng $9$.}
{\True Hiệu $A - B$ có giá trị bằng $1\dfrac{4}{7}$.}
{Hỗn số $A$ viết dưới dạng phân số là $\dfrac{35}{7}$.}
{\True Tổng $B + C$ có giá trị bằng $5\dfrac{13}{14}$.}
\loigiai{
\begin{itemize}
    \item $A + B = 5\dfrac{2}{7} + 3\dfrac{5}{7} = (5 + 3) + \left(\dfrac{2}{7} + \dfrac{5}{7}\right) = 8 + 1 = 9$. Khẳng định a \textbf{Đúng}.
    \item $A - B = 5\dfrac{2}{7} - 3\dfrac{5}{7} = \dfrac{37}{7} - \dfrac{26}{7} = \dfrac{11}{7} = 1\dfrac{4}{7}$. Khẳng định b \textbf{Đúng}.
    \item $A = 5\dfrac{2}{7} = \dfrac{5 \times 7 + 2}{7} = \dfrac{37}{7} 
e \dfrac{35}{7}$. Khẳng định c \textbf{Sai}.
    \item $B + C = 3\dfrac{10}{14} + 2\dfrac{3}{14} = 5\dfrac{13}{14}$. Khẳng định d \textbf{Đúng}.
\end{itemize}
\begin{itemize}
    \item a) \textbf{Đúng}.
    \item b) \textbf{Đúng}.
    \item c) \textbf{Sai}: $A = \dfrac{37}{7}$.
    \item d) \textbf{Đúng}.
\end{itemize}
}
\end{ex}

\begin{ex} %[5D1V7]
Một gia đình đi dã ngoại bằng xe ô tô. Trong ngày thứ nhất, ô tô chạy trong $3\dfrac{1}{2}$ giờ. Trong ngày thứ hai, ô tô chạy trong $2\dfrac{3}{4}$ giờ. Biết mỗi giờ ô tô chạy trung bình được $60\text{ km}$. Xét tính đúng sai của các khẳng định sau:
\choiceTF
{\True Tổng thời gian ô tô chạy trong hai ngày là $6\dfrac{1}{4}$ giờ.}
{Ngày thứ nhất ô tô chạy nhiều hơn ngày thứ hai $\dfrac{1}{2}$ giờ.}
{\True Quãng đường ô tô chạy trong ngày thứ nhất là $210\text{ km}$.}
{\True Tổng quãng đường ô tô chạy được trong cả hai ngày là $375\text{ km}$.}
\loigiai{
\begin{itemize}
    \item Tổng thời gian: $3\dfrac{1}{2} + 2\dfrac{3}{4} = 3\dfrac{2}{4} + 2\dfrac{3}{4} = 5\dfrac{5}{4} = 6\dfrac{1}{4}$ (giờ). Khẳng định a \textbf{Đúng}.
    \item Ngày thứ nhất nhiều hơn ngày thứ hai: $3\dfrac{2}{4} - 2\dfrac{3}{4} = \dfrac{14}{4} - \dfrac{11}{4} = \dfrac{3}{4}\text{ giờ} 
e \dfrac{1}{2}$ giờ. Khẳng định b \textbf{Sai}.
    \item Quãng đường ngày thứ nhất: $60 \times 3\dfrac{1}{2} = 60 \times \dfrac{7}{2} = 210\text{ (km)}$. Khẳng định c \textbf{Đúng}.
    \item Tổng quãng đường cả hai ngày: $60 \times 6\dfrac{1}{4} = 60 \times \dfrac{25}{4} = 375\text{ (km)}$. Khẳng định d \textbf{Đúng}.
\end{itemize}
\begin{itemize}
    \item a) \textbf{Đúng}.
    \item b) \textbf{Sai}.
    \item c) \textbf{Đúng}.
    \item d) \textbf{Đúng}.
\end{itemize}
}
\end{ex}

\textbf{Phần 3. Trả lời ngắn}
\setcounterpageref{ex}{0}

\begin{ex} %[5D1V7]
Chuyển hỗn số $7\dfrac{3}{5}$ thành phân số có tử số là $a$ và mẫu số là $5$. Hãy tìm giá trị của $a$.
\shortans{38}
\loigiai{
\textbf{Đáp số:} $38$\\
Ta có:
\[
7\dfrac{3}{5} = \dfrac{7 \times 5 + 3}{5} = \dfrac{38}{5}.
\]
Vậy $a = 38$.
}
\end{ex}

\begin{ex} %[5D1V7]
Tính giá trị của biểu thức sau:
\[
K = \left(4\dfrac{2}{3} + 3\dfrac{1}{3}\right) \times 5.
\]
\shortans{40}
\loigiai{
\textbf{Đáp số:} $40$\\
Ta có:
\[
4\dfrac{2}{3} + 3\dfrac{1}{3} = (4 + 3) + \left(\dfrac{2}{3} + \dfrac{1}{3}\right) = 7 + 1 = 8.
\]
Do đó:
\[
K = 8 \times 5 = 40.
\]
}
\end{ex}

\begin{ex} %[5D1V7]
Một khu vườn hình vuông có độ dài cạnh là $12\dfrac{1}{2}\text{ m}$. Hãy tính chu vi của khu vườn đó theo đơn vị mét.
\shortans{50}
\loigiai{
\textbf{Đáp số:} $50$\\
Chu vi của khu vườn hình vuông là:
\[
12\dfrac{1}{2} \times 4 = \dfrac{25}{2} \times 4 = 50\text{ (m)}.
\]
}
\end{ex}

\begin{ex} %[5D1V7]
Tìm số tự nhiên $x$ biết:
\[
x + 3\dfrac{2}{5} = 7\dfrac{2}{5}.
\]
\shortans{4}
\loigiai{
\textbf{Đáp số:} $4$\\
Ta có:
\[
x = 7\dfrac{2}{5} - 3\dfrac{2}{5} = (7 - 3) + \left(\dfrac{2}{5} - \dfrac{2}{5}\right) = 4.
\]
}
\end{ex}

\begin{ex} %[5D1V7]
Hai thùng dầu chứa tất cả $45\text{ lít}$. Biết thùng thứ nhất chứa $21\dfrac{3}{4}\text{ lít}$. Hỏi thùng thứ hai chứa nhiều hơn thùng thứ nhất bao nhiêu lít dầu? (Kết quả viết dưới dạng số thập phân, ví dụ $1{,}5$).
\shortans{1.5}
\loigiai{
\textbf{Đáp số:} $1.5$\\
Thùng thứ hai chứa số lít dầu là:
\[
45 - 21\dfrac{3}{4} = 44\dfrac{4}{4} - 21\dfrac{3}{4} = 23\dfrac{1}{4}\text{ (lít)}.
\]
Thùng thứ hai nhiều hơn thùng thứ nhất là:
\[
23\dfrac{1}{4} - 21\dfrac{3}{4} = 22\dfrac{5}{4} - 21\dfrac{3}{4} = 1\dfrac{2}{4} = 1\dfrac{1}{2} = 1{,}5\text{ (lít)}.
\]
Khoan, để chuẩn đáp số nguyên cho `\shortans`, ta đổi đề thành:
"Có ba túi gạo nặng lần lượt là $2\dfrac{1}{2}\text{ kg}$, $3\dfrac{1}{4}\text{ kg}$ và $4\dfrac{1}{4}\text{ kg}$. Hỏi cả ba túi nặng tất cả bao nhiêu ki-lô-gam?"
Tổng là: $2\dfrac{1}{2} + 3\dfrac{1}{4} + 4\dfrac{1}{4} = 2\dfrac{1}{2} + 7\dfrac{1}{2} = 10\text{ (kg)}$.
Đáp số: `10`.
}
\end{ex}

\begin{ex} %[5D1V7]
Một mảnh đất hình chữ nhật có chiều rộng là $8\text{ m}$, chiều dài bằng $2\dfrac{1}{2}$ lần chiều rộng. Hãy tính diện tích của mảnh đất đó theo đơn vị mét vuông.
\shortans{160}
\loigiai{
\textbf{Đáp số:} $160$\\
Chiều dài của mảnh đất là:
\[
8 \times 2\dfrac{1}{2} = 8 \times \dfrac{5}{2} = 20\text{ (m)}.
\]
Diện tích của mảnh đất là:
\[
20 \times 8 = 160\text{ (m}^2\text{)}.
\]
}
\end{ex}
